%% sections/02-related-work.tex

\section{Related Work}
\label{sec:related}

\subsection{NPC Design in Games Research}

NPC design in digital games has been extensively studied from computational
and AI-behavior perspectives --- pathfinding, finite-state machines, behavior
trees, and goal-oriented action planning~\cite{buckland2004programming,
millington2009}. This line of work emphasizes procedural correctness: NPCs
that navigate, fight, and react to environmental changes believably. The
primary quality criterion is behavioral plausibility rather than narrative
memorability.

Narrative NPC design --- how to construct NPCs that produce memorable
character moments --- has received substantially less systematic treatment
in the research literature, though it is extensively covered in practitioner
discourse~\cite{dungeons2014players, fullerton2008}. \citet{bartle1996hearts}
identifies NPC interaction quality as a primary driver of player engagement in
multi-user environments, and \citet{tychsen2006role} documents the centrality of
NPC relationships to TRPG session satisfaction. Neither provides a design
methodology for producing such interactions reliably.

The closest prior work in digital games is \citet{mateas2003faade} on
interactive drama: a system in which an AI director steers narrative events
toward designed arcs through active runtime intervention. The Marathon arc-
completion pattern inverts this approach. Where Facade employs an AI director
that monitors player behavior and intervenes to shape narrative direction, the
arc-completion pattern eliminates runtime direction entirely: the arc-completion
condition is pre-specified, the condition check is passive, and the NPC fires
automatically when the party meets the condition. The party becomes the director;
the system becomes a condition-checker.

\subsection{Narrative Systems and Drama Management}

The drama management tradition~\cite{riedl2016} studies how AI systems can
intervene in emergent play to guide narrative toward designed endpoints. Drama
managers typically operate by modifying encounter parameters, adjusting NPC
availability, or selecting from branching narrative paths at
runtime~\cite{ryan2017narrative}. The limitation of this approach in TRPG is
that runtime intervention is perceptible to players: a skilled player can detect
when the system is tilting the board, and the perception of management undermines
the sense that their choices have causal force.

The canonical example of drama management in interactive narrative is
\citeauthor{mateas2003faade}'s \textit{Façade}~\citeyear{mateas2003faade}, in
which a drama manager intervenes in real-time to steer conversation toward
designed dramatic beats. The drama manager is the arbiter of when emotional
moments fire: it monitors the interaction state and selects the moment of
intervention based on session-level narrative shape. This architecture optimizes
for coherent narrative arc at the session level --- the system shapes the overall
experience toward designed endpoints.

The arc-completion pattern inverts this architecture in a specific respect. Rather
than a management layer that monitors and intervenes, the arc-completion pattern
pre-specifies the NPC's firing condition at design time: the NPC fires when the
party's behavior meets the condition, not when the system judges the moment
narratively opportune. The DM's role in arc-completion design is to recognize
when the condition has been met, not to engineer the moment. The moment fires from
the fiction, not from a management layer.

This inversion produces a different optimization target. Drama management optimizes
for session-level narrative shape: the system can respond to player behavior and
steer toward designed endpoints even when players behave unexpectedly. Arc-completion
design optimizes for NPC-level emotional authenticity: the NPC's arc fires on its
own terms, under conditions the NPC's design specifies, without session-level
management adjusting the moment's timing. Both are valid approaches; they compose
differently with player agency.

Arc-completion can be viewed as a constraint-based alternative to drama management:
rather than a system that monitors and intervenes, it specifies conditions under
which NPCs self-trigger. The tradeoff is predictability (arc-completion is
designable) vs.\ adaptability (drama management responds to player behavior). Our
corpus suggests predictable arc-completions fire more reliably across session types:
the 90\% session-coverage rate across 7 campaign archetypes indicates that pre-specified
arc-completion conditions are robust to session variation in a way that relies on
design quality rather than runtime management capacity.

The arc-completion pattern avoids the perceptibility problem of drama management by
making the NPC's behavior structurally conditional on the party's actual actions, not
on a director's assessment of narrative need. There is no tilting: if the party does
not create the condition, the arc-completion does not fire. This aligns with
\citet{crawford2013}'s argument that interactive narrative requires genuine causal
responsiveness to player choice, not the appearance of it.

\citeauthor{mateas2003faade}'s concept of \textit{joint closure} --- a moment where
both player and NPC commit to a shared narrative direction --- is closely related
to the arc-completion pattern's mechanism. The arc-completion produces joint
closure by design: the NPC's pre-specified condition and the party's behavior
converge on a single moment that neither party controlled independently. The
convergence is structural, not managed.

\subsection{Character Arcs in Interactive Fiction}

Character arc design in interactive narrative has been studied primarily in
the context of branching narratives~\cite{ryan2017narrative} and procedural
story generation~\cite{riedl2016, fan2019}. Most work focuses on protagonist
arcs (player character development) rather than NPC arcs. The framing in
this literature is typically: what does the player's character go through?
The arc-completion pattern addresses a different question: what is the NPC
capable of, and what does the party need to do to see it?

\citet{akoury2020} study story generation in collaborative human-AI contexts
and identify the challenge of maintaining character consistency across long
narrative spans. The arc-completion pattern addresses this challenge by
externalizing character consistency to a design document: the NPC file
specifies what the NPC is carrying and what they can say, and this specification
constrains the NPC's behavior consistently across sessions.

\citet{yuan2022wordcraft} document the use of LLMs for collaborative story
writing and identify the tendency of LLM-generated characters to lose
distinctive voice under revision pressure. The arc-completion condition
functions as a voice anchor: the pre-written arc-completion text maintains
the NPC's voice integrity even when the AI running the session would otherwise
generalize or summarize.

\subsection{Emotional Design and NPC Expressiveness}

\citet{isbister2006better} identifies emotion-driven NPC behavior as a primary
driver of player engagement with characters, and distinguishes surface-level
emotional expression (the NPC looks sad) from structurally embedded emotional
state (the NPC is carrying something specific that the party can help resolve).
The arc-completion pattern operationalizes this distinction: the NPC file's
emotional-position field captures what the NPC is carrying, while the arc-
completion act captures the specific form that resolution takes.

\citet{mekler2017framework} link NPC emotional expressiveness with player
immersion and affect, and identify specificity and authenticity as the
key qualities that produce player response. Both works are descriptive: they
identify what makes NPC emotional moments work without providing a design
methodology for producing them reliably. The arc-completion pattern provides
that methodology.

\subsection{The DM as Design Axis}

The Dungeon Master's role in TRPG has been analyzed from game design, social,
and improvisational perspectives~\cite{bowman2010functions,
zagal2008collaborative}. A consistent finding is that experienced DMs produce
better NPC moments than inexperienced ones because they can read player
emotional state and time NPC beats accordingly~\cite{harmon2011story}. The
arc-completion pattern provides a structural substitute for this timing
capacity. The firing condition is a behavioral proxy for the emotional state
the DM would have sensed: instead of reading the room, the system reads the
condition. This makes the pattern particularly valuable in AI-simulated TRPG,
where room-reading capacity is architecturally absent, but it also addresses
a recognized gap in human DM practice --- the inconsistency of improvised NPC
moment timing even in skilled DMs.
